\begin{tikzpicture}[x=0.70cm,y=22.0cm,
  axis/.style={-{Latex[length=1.4mm]},semithick},
  guide/.style={thin,black!45,densely dotted},
  lab/.style={font=\scriptsize,align=center}]
\draw[axis] (-6.35,0) -- (10.65,0) node[below,font=\scriptsize] {$(V-U_j)/w_j$};
\draw[axis] (-6.35,0) -- (-6.35,0.282) node[above,font=\scriptsize] {$\dd Q/\dd V$ [arb.]};
\foreach \x in {-6,-4,-2,0,2,4,6,8,10}{\draw (\x,0.004) -- (\x,-0.004) node[below,font=\tiny] {\x};}
\foreach \y/\txt in {0.125/0.125,0.250/0.250}{\draw (-6.27,\y) -- (-6.43,\y) node[left,font=\tiny] {\txt};}
\draw[guide] (-6.0,0.25) -- (10.1,0.25);
\draw[-{Latex[length=1.6mm]},very thick] (0,0) -- (0,0.265);
\node[lab] at (-3.1,0.255) {0. Maxwell\\delta limit};
\draw[black!60,semithick] plot[smooth] coordinates {(-6.0000,0.0025) (-5.0000,0.0066) (-4.0000,0.0177) (-3.0000,0.0452) (-2.0000,0.1050) (-1.0000,0.1966) (-0.5000,0.2350) (0.0000,0.2500) (0.5000,0.2350) (1.0000,0.1966) (1.5000,0.1491) (2.0000,0.1050) (2.5000,0.0701) (3.0000,0.0452) (4.0000,0.0177) (5.0000,0.0066) (6.0000,0.0025) (7.0000,0.0009) (8.0000,0.0003) (9.0000,0.0001) (10.0000,0.0000)};
\draw[densely dashed,semithick] plot[smooth] coordinates {(-6.0000,0.0140) (-5.0000,0.0252) (-4.0000,0.0438) (-3.0000,0.0721) (-2.0000,0.1082) (-1.0000,0.1419) (-0.5000,0.1525) (0.0000,0.1562) (0.5000,0.1525) (1.0000,0.1419) (1.5000,0.1264) (2.0000,0.1082) (2.5000,0.0895) (3.0000,0.0721) (4.0000,0.0438) (5.0000,0.0252) (6.0000,0.0140) (7.0000,0.0077) (8.0000,0.0042) (9.0000,0.0022) (10.0000,0.0012)};
\draw[thick] plot[smooth] coordinates {(-6.0000,0.0074) (-5.0000,0.0134) (-4.0000,0.0238) (-3.0000,0.0408) (-2.0000,0.0656) (-1.0000,0.0959) (-0.5000,0.1106) (0.0000,0.1232) (0.5000,0.1322) (1.0000,0.1365) (1.5000,0.1358) (2.0000,0.1305) (2.5000,0.1213) (3.0000,0.1097) (4.0000,0.0833) (5.0000,0.0587) (6.0000,0.0392) (7.0000,0.0251) (8.0000,0.0156) (9.0000,0.0094) (10.0000,0.0056)};
\node[lab,black!60] at (2.75,0.212) {1. intrinsic\\$w_j=n_jRT/F$};
\node[lab] at (-4.15,0.123) {2. ensemble\\$\sigma_\eta=1.25\sigma_\mathrm{int}$\\scale $1.6w_j$};
\draw[-{Latex[length=1.2mm]},gray] (4.2,0.094) -- (7.8,0.030);
\node[lab] at (7.15,0.101) {3. one-sided tail\\$L_V=1.5w_j$};
\draw[{Latex[length=1.1mm]}-{Latex[length=1.1mm]},thin] (-1.6,0.018) -- (1.6,0.018)
  node[midway,above,font=\scriptsize] {$1.6w_j$};
\draw[{Latex[length=1.1mm]}-{Latex[length=1.1mm]},thin] (1.0,0.050) -- (2.5,0.050)
  node[midway,above,font=\scriptsize] {$L_V$};
\end{tikzpicture}
\caption{두-상 broadening 의 폭 예산(식~\eqref{eq:widthbudget}) 개선안. 좌표는 T8\_calc.py 로 식 그대로 수치 평가했다.
평형 델타에 ② 내재 logistic 미분 종(폭 $w_j$, 회색 실선)이 얹히고, ③ 앙상블 $\eta$ 산포와의 합성곱은
$\sigma_\eta=1.25\,\sigma_\mathrm{int}$ 예시에서 유효 scale $1.6w_j$ 의 대칭 종(파선)으로 넓어진다.
여기에 ① 유한율속 단측 지수 커널($L_V=1.5w_j$)을 수치 적분으로 합성하면 진행 방향 꼬리(굵은 실선)가 생긴다.
대칭 두 몫은 $w_j$ 로 흡수되고 비대칭 꼬리는 $L_{V,j}$ 로 분리된다.}
\label{fig:widthbudget}
